% !TeX root = ../main.tex
\section*{How to read this document}
\addcontentsline{toc}{section}{How to read this document}

\paragraph{Purpose.} Coding agents work on the transcription pipeline of Guitar
Virtual Amp, run benchmarks and change models. This document is the place where
their work is explained in one coherent story: what problem each piece solves,
what the literature knows about it, what the app does, how well it works, and
what is planned. It replaces an earlier review written in conference-paper
style. It is deliberately longer and more explicit than a paper.

\paragraph{Structure.}
\begin{itemize}
\item \textbf{Part I, Background} (Sections~\ref{sec:problem}--\ref{sec:eval})
  defines the task and teaches the minimum of audio processing, machine
  learning and evaluation needed to follow the rest. Skip it if you know these.
\item \textbf{Part II, The pipeline} (Sections~\ref{sec:datasets}--\ref{sec:e2e})
  walks through the stages in the order audio flows through them. Each section
  follows the same template: \emph{what the stage must do}, \emph{how methods
  work}, \emph{state of the art}, an ``In Guitar Virtual Amp'' box, and
  \emph{failure modes}.
\item \textbf{Part III, Our system} (Sections~\ref{sec:app}--\ref{sec:assessment})
  shows the whole app pipeline, all our benchmark results with their protocol,
  and an assessment: where we stand against the state of the art, which
  decisions were taken and why, and what the agents are working on next.
\item \textbf{Appendices}: a glossary and a map from source files to sections,
  with the commands that reproduce each benchmark.
\end{itemize}

\paragraph{Labels.} The coloured boxes have fixed meanings:
\begin{keyidea}[]
The intuition behind a method, before its details.
\end{keyidea}
\begin{inapp}[]
What the app does, with file names under \file{guitar_core/src/} unless another
path is given, and the status of each part.
\end{inapp}
\begin{caution}[]
A pitfall: a number that is easy to misread, a comparison that is not fair, a
known gap.
\end{caution}
Inline tags say where a claim comes from: \measured{} means we ran it in this
repository (the script is named), \reported{} means a paper states it (with a
citation), and \expected{} means it is a hypothesis that no one has measured
yet. For parts of the app: \done{} exists in the code, \wip{} is being worked
on, \planned{} is only in the roadmap.

\paragraph{Related files.} \file{ROADMAP.md} (project status),
\file{docs/transcription_roadmap.md} (the detailed research log, from which most
of the measured numbers here are taken), \file{CLAUDE.md} (architecture notes for
agents), \file{scripts/train/README.md} (training and Python benchmarks). The
companion document on estimating a recording's guitar tone is
\file{docs/papers/tone-estimation/}.

\paragraph{Snapshot.} Numbers and code descriptions are as of 25 November 2026
(commit \code{25afb6d} plus the uncommitted work described where relevant).
When the code changes, the ``In Guitar Virtual Amp'' boxes and Part~III are the
parts to update.

\subsection*{The whole story in one table}
\addcontentsline{toc}{subsection}{The whole story in one table}

Table~\ref{tab:summary} summarises the rest of the document: for each stage,
the best published approach, what the app uses, our best measured number and
the main remaining gap. Every entry is explained later.

\begin{table}[h]
\centering
\caption{Stage-by-stage summary. F50 = note F1 with a 50\ms{} onset window
(Section~\ref{sec:eval}). ``v2'' = our electric-guitar fine-tune of the GAPS
model. All numbers in the last two columns are \measured{} unless marked.}
\label{tab:summary}
\small
\begin{tabularx}{\textwidth}{@{}>{\raggedright}p{2.0cm}>{\raggedright}p{3.1cm}>{\raggedright}p{3.0cm}>{\raggedright}p{3.4cm}X@{}}
\toprule
Stage & Best published & App uses & Our best number & Main gap \\
\midrule
Separation (\S\ref{sec:separation}) & Band-split transformers (BS-RoFormer) & BS-RoFormer SW, else Demucs \code{htdemucs_6s} & Real songs: F50 0.41 (mix) $\to$ 0.70 (stem) & All guitars in one stem \\
Note transcription (\S\ref{sec:notes}) & Guitar CRNNs (GAPS); MuScriptor on mixes & GAPS CRNN; v2 as local option & EGSet12 (real amp) 0.89; GOAT amp renders 0.71 & Distorted guitar in real mixes: 0.56 \\
Rhythm (\S\ref{sec:rhythm}) & Beat This! & Beat This! final0 + bar Viterbi & Beat F 0.97 (accompaniment), 0.55 (solo) & Solo lines; compound meters \\
Fingering (\S\ref{sec:fingering}) & Learned seq2seq, audio-conditioned (TART) & Learned tagger (local), else Viterbi & String accuracy, true notes: 0.62 solo, 0.81 EGDB & Few players in training data \\
Techniques (\S\ref{sec:techniques}) & Learned per-note classifier (TART) & Rules on the pitch contour & Bends found 35--56\%, vibrato 20--40\% & No palm mute, harmonics; distortion \\
Whole tab (\S\ref{sec:benchmark}) & TART: 54\% Tab F1, solo guitar \reported{} & All of the above & Tab F1 0.65 on separated mixes (tagger) & No real-song benchmark yet \\
\bottomrule
\end{tabularx}
\end{table}
